\section{Por qué la generación y la comprensión no se han fusionado realmente}

La sección anterior explicó la diferencia de fondo de las imágenes estáticas; esta sección examina los tropiezos concretos de Zhang Xiangyu (张祥雨) con los sistemas multimodales. Él intentó reunir en un mismo sistema la generación de imágenes, la comprensión de imágenes, la generación de texto y la comprensión de texto: tokenizar las imágenes, procesar lo multimodal con un Transformer de estilo lingüístico y conectar después módulos de generación como diffusion. El resultado fue que el modelo de comprensión se volvió cada vez más capaz y el modelo de generación también, pero al juntarlos no apareció un efecto acumulado en el que 1+1 fuera mayor que 2.

La evidencia más contundente es esta: al quitar la parte de generación, la capacidad de comprensión no se ve afectada; al conectar el modelo de comprensión al modelo de generación, la controlabilidad de la generación sigue siendo deficiente. El modelo puede reconocer que un video contradice el sentido común físico, pero no puede impedirse a sí mismo generar videos que lo contradicen. Esto indica que las dos ramas todavía no comparten un mismo proceso interno de razonamiento.

\begin{figure}[H]
\centering
\includegraphics[width=0.94\textwidth]{generation-understanding-gap.png}
\caption{Generación y comprensión sin fusionar: ambas ramas se fortalecen, pero juntas no dan 1+1 mayor que 2. Diagrama conceptual propio, elaborado a partir de la conversación entre 00:29:10 y 00:38:45.}
\end{figure}

\begin{knowledgebox}{Lectura de la figura: conectar sistemas no equivale a fusionar capacidades}
A la izquierda, el modelo de comprensión destaca en ver imágenes y videos y en alinearlos con el lenguaje; a la derecha, el modelo de generación destaca en producir imágenes y videos. El gap del centro no es una interfaz de ingeniería, sino la imposibilidad de que las capacidades se corrijan mutuamente. Una integración verdadera debería hacer que la comprensión restrinja la generación y que el proceso de generación, a su vez, entrene la comprensión.
\end{knowledgebox}

\subsection{El problema de fondo de la controlabilidad de la generación}

Esta sección desglosa la «baja controlabilidad». Los modelos de generación de imágenes y video pueden aprender patrones visuales de alta probabilidad, pero las escenas complejas exigen respetar restricciones geométricas, físicas, de consistencia de identidad, de continuidad temporal y semánticas. Las indicaciones en lenguaje no siempre logran transmitir con precisión estas restricciones al proceso de generación; tampoco el modelo de comprensión logra necesariamente convertir el «esto no es razonable» en una corrección ejecutable dentro del modelo de generación.

\begin{knowledgebox}{Términos clave: generación, comprensión, alineación, controlabilidad}
\noindent\begin{tabularx}{\textwidth}{p{0.18\textwidth}p{0.36\textwidth}X}
\toprule
Concepto & Significado & Problema en este episodio \\
\midrule
Generación & Producir imágenes, videos o texto & Puede ser realista, pero no necesariamente controlable. \\
Comprensión & Reconocer objetos, relaciones, acciones y semántica & Puede señalar errores, pero no siempre guía la generación. \\
Alineación & Coincidir con la intención y la evaluación humanas & La alineación de imágenes depende del lenguaje, las preferencias y la definición de la tarea. \\
Controlabilidad & Generar de forma estable el resultado buscado bajo restricciones & Las restricciones de múltiples objetos, físicas y temporales son las que más fallan. \\
\bottomrule
\end{tabularx}
\end{knowledgebox}

\begin{warningbox}{El modelo de generación «sabe que está mal», pero aun así lo hace mal}
No es una contradicción superficial. Juzgar si un resultado es erróneo y evitar los errores paso a paso durante la generación son dos capacidades distintas. La primera se parece a calificar; la segunda, a planificar y ejecutar. Lo que les falta a los sistemas multimodales es justamente la cadena de razonamiento que conecte la calificación con la ejecución.
\end{warningbox}

\subsection{Resumen de la sección}

La dificultad de la integración multimodal radica en que la comprensión y la generación no se refuerzan mutuamente con solo conectarlas. Necesitan compartir un proceso intermedio que pueda explicarse, explorarse mediante búsqueda y corregirse; esa es precisamente la razón por la que más adelante entran en la discusión o1, CoT y el Long CoT visual.
